\documentclass[12pt,a4paper]{article}

% Idioma e codificação
\usepackage[T1]{fontenc}
\usepackage[utf8]{inputenc} % Overleaf já usa UTF-8, mas não faz mal manter
\usepackage[portuguese]{babel}

% Pacotes úteis
\usepackage{lmodern}    % fontes
\usepackage{microtype}  % melhor tipografia
\usepackage{hyperref}   % URLs e links clicáveis
\usepackage{geometry}
\geometry{margin=2.5cm}

% Título
\title{Método de Avaliação de Usabilidade (SUS)}
\author{Pedro Tales / InfoSuplmentos}
\date{\today}

\begin{document}

\maketitle
\tableofcontents
\bigskip

\section{Método de Avaliação de Usabilidade}

\subsection{Objetivo}
O objetivo desta etapa é avaliar a usabilidade e a experiência geral de uso do site de suplementos alimentares, considerando aspectos de navegação, clareza da informação, estética e satisfação geral dos usuários. A avaliação foi conduzida por meio de um questionário desenvolvido no Google Forms, baseado em princípios do System Usability Scale (SUS) e adaptado às características do projeto.

\subsection{Participantes}
Participaram da avaliação usuários potenciais do site e estudantes de turma. Foram coletadas informações demográficas como idade, escolaridade e frequência de compra de suplementos, permitindo caracterizar o perfil do público-alvo.

\subsection{Instrumento de Coleta de Dados}
O instrumento utilizado foi um formulário eletrônico elaborado no Google Forms (\url{https://docs.google.com/forms/d/e/1FAIpQLSedlxPpPGPA_pH3LVlV916ewnRjc4-TeTUIGmWmLySkYKcMcA/viewform}), composto por quatro seções:

\begin{enumerate}
  \item \textbf{Perfil do Usuário}: coleta de informações como idade, gênero, escolaridade e frequência de consumo de suplementos.
  \item \textbf{Acesso e Navegação}: questões sobre o dispositivo de acesso, facilidade de navegação, clareza dos menus e organização do conteúdo.
  \item \textbf{Conteúdo e Informações}: avaliação da clareza das descrições dos produtos, utilidade das imagens e confiança nas informações apresentadas.
  \item \textbf{Satisfação e Sugestões}: perguntas sobre o nível geral de satisfação, probabilidade de recomendar o site e campo aberto para sugestões.
\end{enumerate}

\subsection{Procedimento}
Os participantes foram convidados a acessar o site de suplementos e explorar suas funcionalidades principais (busca de produtos, visualização de detalhes e leitura de informações nutricionais). Após a navegação, responderam ao questionário de avaliação. As respostas foram registradas automaticamente no Google Sheets, exportadas em formato CSV e analisadas quantitativamente.

\subsection{Análise de Dados}
A análise quantitativa envolveu:
\begin{itemize}
  \item Cálculo da pontuação SUS individual e média (0–100).
  \item Estatísticas descritivas: média, desvio-padrão, mediana, e intervalo interquartil.
  \item Distribuições de frequência por categoria (perfil, dispositivo de acesso, satisfação geral).
  \item Representações gráficas: histogramas, boxplots e gráficos de barras.
\end{itemize}

\subsection{Ética e Privacidade}
Os participantes responderam de forma anônima, sem coleta de dados pessoais sensíveis. As respostas foram utilizadas exclusivamente para fins acadêmicos e de melhoria do projeto.

\end{document}
